problem:  
Let \( (M^4,g) \) be a compact, simply connected Riemannian 4‑manifold with nonnegative sectional curvature admitting an effective, isometric \(S^1\)-action. Denote the quotient by \(M^* = M / S^1\), which is an Alexandrov 3‑space with curvature \(\ge 0\). Its singular locus \(\Sigma^*\subset M^*\) is the image of orbits with nontrivial isotropy; the set of its 2‑dimensional faces corresponds to fixed‑point 2‑spheres and its 1‑dimensional edges correspond to exceptional orbits.  

The local geometry near each stratum is well understood: the slice theorem implies that each boundary component of \(M^*\) has a local product model, and each edge has a local cone model on an \(S^2\) with a circle of cone angle \(2\pi/m\) determined by the isotropy weight \(m\).  However, it is not known whether and how these local cone models globally constrain the **global Alexandrov curvature** and inter‑angle relations along \(\Sigma^*\).  

**Problem:**  
Assume \(M\) is simply connected, \(\mathrm{sec}\ge0\), and \(S^1\) acts isometrically.  Let \(\Sigma^*\subset M^*\) be the graph of singular strata with associated isotropy weights \(\{m_e\}\) along the edges.  
> Are there global **Alexandrov curvature compatibility inequalities** among these integers \(\{m_e\}\), forced by \(\mathrm{curv}(M^*)\ge 0\)?  

Equivalently:  
> Can the nonnegative curvature of \(M\) impose metric constraints on the way cone singularities (arising from different isotropy weights) meet at vertices of \(\Sigma^*\) in \(M^*\)?  

Explicitly, one may ask:
- Are there inequalities bounding the sum of cone angles around a vertex (encoded by the isotropy integers on adjacent edges) required by \(\mathrm{curv}(M^*)\ge0\)?  
- Do Alexandrov angle sum conditions—analogous to Gauss–Bonnet for polyhedral surfaces—create new restrictions on the possible configurations of singular edges meeting at the images of fixed‑point components?

---

Why is it a "good" problem:  
This question moves beyond local slice geometry (which is representation‑theoretically rigid) to global Alexandrov curvature compatibility—something currently unknown.  It asks whether **curvature bounds propagate through discrete symmetry data** (the isotropy weights and their combinatorics), yielding inequalities that might tighten or even classify the possible singular graphs in orbit spaces of nonnegatively curved \(S^1\)‑manifolds.

Its significance lies in several aspects:

- **New bridge between discrete and continuous geometry:**  The problem connects integral isotropy data (\(m_e\)) from transformation group theory with Alexandrov curvature bounds, inviting a hybrid of metric comparison and combinatorial analysis.
- **Potential classification advancement:**  If such inequalities exist, they would further restrict the topologically possible Fintushel weight configurations realizable by metrics with \(\mathrm{sec}\ge0\).
- **Analogy with polyhedral geometry:**  The question echoes the classical angle‑sum constraints (e.g., in Euclidean vs. spherical polyhedra), but now in the context of 3‑dimensional Alexandrov spaces arising from smooth 4‑manifolds—completely unexplored terrain.
- **Conceptually simple, profoundly deep:**  “Do curvature ≥ 0 metrics constrain how isotropy cones fit together in the quotient?”  The statement is elementary but resolving it would link Alexandrov comparison geometry, Riemannian group actions, and 4‑manifold topology in a fundamentally new way.

Thus, it epitomizes the “narrow entrance, profound depth’’ criterion: a single clarity‑laden question whose answer could reveal an uncharted rigidity principle governing how global curvature interacts with the discrete symmetry structure of 4‑manifolds.